\chapter{Población y muestra}
\label{cap:poblacion}

\section{Población}

La población estudiada tiene tres grupos. El primero es el personal de la
agencia que participa en el proceso de reservas: administrador, encargados
de atención o ventas y personal de coordinación, si corresponde. El segundo
son los clientes que soliciten información o realicen reservas durante el
periodo de estudio. El tercero son las reservas registradas en ese mismo
periodo, que se revisarán como fuente documental. Los proveedores serán
consultados solo como informantes complementarios.

La cantidad exacta de personas y reservas no puede definirse todavía,
porque falta seleccionar la agencia y conocer sus registros. Por ello no se
utilizan cifras supuestas, y los totales se anotarán una vez identificada
la institución.

\section{Muestra}

Para el personal relacionado con las reservas se propone trabajar con todos
los participantes del proceso, siempre que el número y la autorización lo
permitan. Para los clientes se plantea, inicialmente, entrevistar hasta 20
personas mediante selección por conveniencia, considerando a quienes
consulten o reserven durante el periodo de estudio. Esta cantidad es una
meta de trabajo y no representa una muestra estadística definitiva.

Si se necesita una muestra probabilística, será necesario conocer
previamente la cantidad total de clientes y establecer el nivel de
confianza y margen de error. La decisión final se tomará según los datos
disponibles y el propósito de la evaluación.

\begin{tablaapa}{tab:poblacion}{Plan preliminar de población y muestra}
\begin{tabular}{@{}p{4.6cm}p{5.2cm}p{5.2cm}@{}}
\toprule
Grupo de estudio & Población & Muestra prevista \\
\midrule
Personal responsable de reservas
  & Total por confirmar con la agencia
  & Censo de quienes participan en el proceso \\
Clientes que consultan o reservan
  & Total por confirmar en el periodo de estudio
  & Hasta 20 clientes, propuesta por conveniencia \\
\bottomrule
\end{tabular}
\nota{Elaboración propia.}
\end{tablaapa}
